% =====================================================================================================
% Introduction (no heading, AER style). Version 3. Writer: section writer "introduction" (round 3), 2026-09-24.
% v2 text saved in paper/sections/v2/introduction.tex. Numbers verified against output/tables/*.csv (rb1_sigmaC,
% rb2_decisions, rb3_econ2, ra2_wedge); sources, referee comments and open issues are listed in
% paper/notes/writer3_introduction.md, which also holds the abstract, JEL codes and keywords for main.tex.
% [rb4 integration, 2026-09-24] sigma* numbers after Chinchilla's FLOP rebuild: paper/notes/integration_log_rb4.md.
% [round 3, WP8, 2026-09-25] Items H3-H10 of paper/notes/round3_fixlist.md: virtual value of compactness, sigma* by
% compute, 49 budget-level decisions, the rebuilt sign result (rb5_sign), the trend with a developer bootstrap,
% "recipe-dependent" extrapolation (decision D-5), the general-isocost qualifier. Every number and its source file:
% paper/notes/round3_WP8_log.md. The red TBD-m9 block follows Amendment 2, item 10, and is filled by the sentences of
% Amendment 2, item 4 ("Placeholder sentences of the introduction and the conclusion").
% =====================================================================================================

Training a frontier language model has become one of the largest single investments in the technology sector: frontier
training compute has grown four- to fivefold a year since 2018 \citep{sevilla2024training}, and its amortized cost
about 2.4-fold a year since 2016 \citep{cottier2024rising}. Labs, our term for both the laboratories that run training
experiments and the developers that release models, plan this work with neural scaling laws, which relate a model's
loss $L$ on held-out text to its parameters $N$ and training tokens $D$ \citep{kaplan2020scaling,
hoffmann2022training}. Training costs about $6ND$ floating-point operations (FLOP) and serving about $2N$ FLOP per
token, so a scaling law is a production function with a known cost side, and the choice of $N$ and $D$ depends on what
else a lab values: one that expects heavy use of a model \citep{sardana2024beyond}, or wants it to run on a laptop
\citep{benallal2025smollm2}, trains a smaller model on more data. Two objects govern this choice, how easily
parameters and data substitute and how much labs value a compact model; both enter forecasts of data demand and of the
returns to compute in models of growth. Labs' optimization cuts both ways for measuring them.

A lab that minimizes cost places its models on the expansion path, where they carry almost no information about the
curvature of the technology. The same optimization makes its choices informative about its objectives: the further a
model lies from the path that minimizes training cost, the more the lab must value something that training cost leaves
out. Designed experiments break the circle: they supply the technology that turns choices into revealed objectives.

\textit{Theory.} Because training compute is multiplicative in parameters and data, isocost lines are straight in
logs, and cost minimization equates the output elasticities of the two inputs, $\varepsilon_N$ and $\varepsilon_D$. An
interior optimum then requires an elasticity of substitution $\sigma$ between zero and one (Lemma~\ref{lemma:sigma});
U-shaped IsoFLOP profiles, loss against model size at fixed compute, are its empirical content. The ratio
$w\equiv\varepsilon_N/\varepsilon_D$ measures what else a lab values. A lab stops adding parameters when their
marginal product equals their full marginal cost, which includes whatever it loses when a model of given quality is
larger; tokens cost only training. Under any objective that rewards lower loss, and absent data costs, caps on tokens
and size-dependent FLOP prices, $w-1$ is therefore the value of compactness at given quality per unit of training
cost, including the shadow value of any binding limit on size (Proposition~\ref{prop:wedge}); we report it as the
compactness share $s\equiv(w-1)/w$. Compactness can pay through serving cost, adoption, memory tiers or latency; the
wedge measures their sum.

On the expansion path, labs' data identify the path and the loss--compute frontier, but the curvature of the
technology only through functional form (Proposition~\ref{prop:ident}). Off the path, information about the
compute-optimal ratio of tokens to parameters, $M^*$, is second order in labs' allocation errors, and information
about curvature is fourth order: curvature affects loss only through the loss from misallocation, which is itself
second order. Designed experiments restore identification without a functional form: with $\sigma^*$ the elasticity on
the expansion path, $1/\sigma^*-1$ equals the curvature of the IsoFLOP profile in log model size at its minimum,
divided by twice the slope of the loss--compute frontier (Proposition~\ref{prop:modelfree}). Beyond the largest
designed budget, the sign of $w-1$ needs even less: for any technology whose IsoFLOP profiles have a single peak and
no other critical point, a model with $M\equiv D/N$ trained with compute $C$ is over-trained exactly when $M>M^*(C)$,
so bounds on $M^*(C)$ identify the sign (Proposition~\ref{prop:pi}).

\textit{The technology.} We apply the estimator to the annealed IsoFLOP experiments of three laboratories, Chinchilla
\citep{hoffmann2022training}, Llama~3 \citep{grattafiori2024llama} and Marin's ladders on three corpora
\citep{marin2026ladders}, and to the near-factorial Farseer sweep \citep{li2025predictableb}. Where the designs
overlap, $10^{19}$ to $3\times10^{20}$ FLOP, $\sigma^*$ is about 0.7 with no common drift; one estimate per study
gives 0.69, with a 95 percent confidence interval of $[0.64, 0.73]$, and FLOP and parameter conventions, estimators
and bandwidths move this mean between 0.66 and 0.70. At $6\times10^{20}$--$3\times10^{21}$ FLOP, in the two IsoFLOP
designs that reach those budgets, $\sigma^*$ is about 0.59 (Farseer's local path gives 0.66 at $10^{21}$), and beyond
the designs it is only partially identified. On Llama~3 and Farseer, parametric fits come close to these values only
when the outer curvature of the loss function is free. The level of the technology does not travel: at $10^{21}$ FLOP,
$M^*$ differs more than tenfold across the sweeps and published laws we use. Our own experiment, with a pre-analysis
plan committed before estimation, trains small transformers at $2\times10^{14}$ to $10^{17}$ FLOP on two corpora that
differ in quality. \textcolor{red}{[TBD-m9: one sentence per Amendments~1 and~2: model-free $\sigma^*$ by corpus, $[x,
y]$; neutrality of the filter; extrapolation slope against its null.]}

\textit{What over-training reveals.} We invert labs' first-order conditions on a clean sample of 77 dense open-weight
base models released in 2023--2025 by developers that disclose training tokens. Family members trained on one token
budget reveal one number, so the models form 49 allocation decisions. In the families we fit, $\ln
w=(1/\sigma^*-1)\ln[M/M^*(C)]$, so the inversion needs the curvature and the zero point $M^*(C)$ at each model's
compute. Under our reference technology, the Chinchilla data fitted with the outer curvature free, the median decision
was made as if a model one percent smaller at the same loss were worth 2.8 percent of its training cost ($s=0.73$),
with a 95 percent interval for $s$ of $[0.51, 0.81]$ under a cluster bootstrap. The level cuts both ways: a zero-point
allowance as large as the shift across DataDecide's recipes lowers the models' median lower bound on $s$ to 0.03, a
curvature that keeps falling with compute raises the median decision's $s$ to 0.87 at $\sigma^*=0.60$ (zero points
fixed), and the 20 decisions under a possible cap on tokens are lower bounds; the 29 point estimates have a median of
0.74. The evidence is more consistent with adoption and memory tiers than with serving cost: models built for devices
are the most over-trained and open-weight models more so than closed ones at given compute, while too few developers
serve their own sizes for a serving test to have power.

\textit{The sign and the trend.} The sign of over-training is more robust than its size. With $M^*$ extrapolated along
the IsoFLOP designs' own paths from their largest bracketed budgets, with the sampling error of level and slope, $w>1$
is identified for 78 percent of models and 69 percent of decisions. The share is 53 percent if labs' $M^*$ may differ
by a factor of 3.4, the shift across data recipes in DataDecide \citep{magnusson2025datadecide}, 13 percent if only
normal inputs bound the path and 23 percent if any technology in our set may anchor it, and it stays above one half up
to a factor of about 3.7. Identified models carry only 24 percent of training compute, because flagships lie close to
the path. Inside the designs, parametric forms understate the wedge at high $M$ in all three recipes that allow the
check, but the local log wedge is convex in two and concave in one, so the direction of extrapolation error in $M$
beyond the designs is recipe-dependent. From 2023 to 2025 revealed over-training rose steeply: under the reference the
compute-weighted share over decisions rose from 0.24 to 0.81 and the median from 0.49 to 0.74, under every technology
in our set and without any single developer, though a bootstrap over developers establishes the rise over the two
years but not the step from 2024 to 2025. The trend largely restates the rise in tokens per parameter, and measures a
rising value of compactness only if developers optimized against technologies like ours.

\textit{Why these objects matter.} Which objects matter is itself a finding (Section~\ref{sec:econ}). The allocation
exponent, among the least portable objects, drives forecasts of data demand: across its estimates, compute-optimal
data demand grows 2.0- to 2.9-fold a year at current compute growth. For the median open-weight decision, the rise in
over-training since 2023 has multiplied tokens per model at given compute by about 2.1 a year (95 percent interval
0.95 to 4.4). The curvature $\sigma^*$ scales the revealed value of compactness but barely moves the cost of a binding
cap on unique data when repeated data lose value slowly; as such caps bind, measured wedges will understate the value
of compactness, so a flattening of the trend would not show that labs value it less. Levels hinge on $M^*$ at frontier
compute, which no public design observes; IsoFLOP profiles published at frontier budgets would turn our signs into
levels.

The methods travel beyond language models. The identification result gives rates to the observation of
\citet{marschak1944random} and \citet{bond2005adjustment} that optimizing firms facing common prices do not reveal
their technology. The model-free estimator, with a correction for the curvature of the isocost itself, applies
wherever experimenters vary two inputs at known cost (Online Appendix Corollary~\ref{cor:A-general}), as in factorial
fertilizer trials, chip-design sweeps and dose-combination trials. The sign result identifies the direction of a
first-order-condition wedge from the location of the cost-minimizing input ratio alone, which bears on measuring
markups and input distortions \citep{deloecker2012markups, raval2023testing} wherever designed experiments locate that
ratio, as in engineering or agronomic trials; production data alone do not.

\textit{Related literature.} \citet{villalobos2023trading}, \citet{erdil2024optimally}, \citet{sardana2024beyond},
\citet{bian2025scaling} and \citet{roberts2026test} solve the forward problem: optimal size and data given inference
demand, latency or repeated sampling at test time. We solve the inverse problem; their mechanisms, the cost of serving
\citep{erdil2025inference} and who bears it determine what the recovered value of compactness contains.
\citet{hao2026theory} model a developer that pays for training and serving under a Leontief approximation, which rules
out over-training; Proposition~\ref{prop:wedge} inverts such problems when $\sigma>0$, also for labs that do not bear
serving costs. \citet{kricheli2026tokens} show that fixed-ratio designs leave scaling-law parameters ill-conditioned
when the two exponents are close; on a compute-optimal path the degeneracy arises whatever the exponents
(Proposition~\ref{prop:ident}), and economic wedges supply the design spread that a curvature parameter needs
\citep[][see also \citealp{ackerberg2015identification, gandhi2020identification}]{kiefer1959optimum,box1959design}.
The inversion follows the production approach of \citet{deloecker2012markups}, but $w$ is not a markup, and its
elasticities come from experiments, which avoids the circularity stressed by \citet{doraszelski2021reexamining} at the
cost of external validity; the critiques of \citet{bond2020unpleasant} and \citet{raval2023testing} carry over
(Section~\ref{sec:wedge}). A companion paper \citep{okar2026observational} takes the framework to observational data
and algorithmic progress \citep{mertens2026secret, whitfill2025note}.
